
% Actual theta constant-term producers, not conditional restatements.
% Requires the actual models/evaluations and the convergence estimates proved
% in the preceding sections. The seminorm consequence needed here is proved
% explicitly below from the same coefficient majorants.

\section{The two polarizations and the theta constant terms}
\label{sec:sr-theta-constant-terms}

All oscillator actions in this section are the common genuine actions
fixed above, with the specified splitting characters and additive
character. The evaluations \(e_0,e_A,e_B\) are the corresponding partial
Fourier and evaluation maps. Thus \(e_0\) removes the plane in \(W\),
whereas \(e_A,e_B\) remove the planes in both \(X\) and \(W\).
No choice of a new local matching quotient is made in this section.
Write \(\Theta(f)\) for the rational theta functional in the indicated
model, and put
\[
 T_A(\xi,f)=\Theta\bigl(e_A\rho(\xi x\otimes y')f\bigr),
 \qquad
 T_B(\xi,f)=\Theta\bigl(e_B\rho(\xi y\otimes x')f\bigr)
 \quad(\xi\in E).
\]
In these two definitions the coefficient is \(\xi\) itself; the
conjugation in \(u_B(l)=\bar l\,y\otimes x'\) will enter when we unfold
the first unitary group.

\begin{srlemma}[Poisson comparison of the two line sums]
\label{lem:sr-theta-two-line-poisson}
There is \(c_{AB}\ne0\), depending on the two evaluations and their
fixed Haar normalizations, such that for every adelic Schwartz vector
\(f\),
\[
 \sum_{\xi\in E}T_A(\xi,f)
        =c_{AB}\sum_{\xi\in E}T_B(\xi,f).
 \tag{T1}\label{eq:sr-theta-two-line-poisson}
\]
Both series are absolutely convergent. The scalar is fixed before the
finite splitting and Hecke character packet.
\end{srlemma}

\begin{proof}
Let
\[
 I_\cap=Ex\otimes Ex'+X_1\otimes Ex'+Ex\otimes W_1.
\]
The two evaluation spaces are
\(I_A=I_\cap+Ey\otimes Ex'\) and
\(I_B=I_\cap+Ex\otimes Ey'\).
Reduce first by \(I_\cap\). The remaining symplectic space is the
orthogonal sum of \(X_1\otimes W_1\) and the plane spanned by
\(x\otimes y'\) and \(y\otimes x'\). Their tensor pairing is a unit,
with the sign forced by
\(\langle y',x'\rangle=-1\).
Consequently the latter summand has one \(E\)-valued Schr\"odinger
coordinate. In these coordinates the two evaluations are, respectively,
evaluation in this coordinate and evaluation after its Fourier
transformation. They have the same reduced Heisenberg action on
\(X_1\otimes W_1\).

To spell out the normalization issue, the constructed partial Fourier
maps identify the common reduced model with
\(\mathcal S(\mathbb A_E)\otimes
 \mathcal S(X_1\otimes W_1)\).
An intertwiner with the stated evaluation invariances is a nonzero
scalar multiple of these coordinate maps: translation invariance in an
integrated coordinate and modulation invariance in an evaluated
coordinate force Haar integration and evaluation at zero, respectively.
Applying this one coordinate at a time proves the assertion, including
for the actual evaluations. The scalars are fixed by the evaluations
and their measures; finite unramified vacuum normalization fixes their
finite local scalars. They do not depend on a splitting character on a
unitary Levi, since the present operators involve only the Heisenberg
action and the fixed rational theta functional.

Thus, after fixed nonzero scalar factors, the two sides of (T1) are
\[
 \sum_{\xi\in E}\sum_{\eta}
        F(\xi,\eta),\qquad
 \sum_{\xi\in E}\sum_{\eta}
        \widehat F(\xi,\eta)
\]
in a common adelic Schwartz model. Here \(\eta\) runs over the rational
coordinates of the lower model, and the hat is Fourier transformation
in the \(E\) coordinate for the actual trace character. A possible
reflection or conjugation of this coordinate permutes \(E\) and has
no effect on the sum. Poisson summation in that coordinate proves (T1).

For completeness, these sums and the uses of Poisson summation are
absolute: finite support and local constancy leave finitely many
cosets of a fractional lattice at infinity, and every infinite
Schwartz seminorm bounds a summand by
\(C_N(1+\|\xi\|+\|\eta\|)^{-N}\).
Lattice packing makes the latter sum finite for \(N\) larger than its
real dimension. The same argument is uniform on compact families of
Heisenberg parameters. Hence all the preceding regroupings and partial
Fourier transformations are justified.
\end{proof}

\begin{srlemma}[The full-rank cutoff under the \(B\) evaluation]
\label{lem:sr-theta-cone-cutoff}
In the specified split local Fourier model the upper cutoff means
support on matrices of full row rank \(n=\dim_E X\).
For a vector \(f\) in this cutoff space,
\[
 e_Bf=0,\qquad
 e_B\rho(a\,y\otimes x')f\in\mathcal S^o(X_1\otimes W_1)
                 \quad(a\in\mathbb A_E).
 \tag{T2}\label{eq:sr-theta-cone-cutoff}
\]
The lower cutoff uses the same split place and the same embedding.
If \(e_0f\) is in the corresponding cutoff for \(X\otimes W_1\), then
\[
                         e_Af=0.                 \tag{T3}
 \label{eq:sr-theta-zero-A}
\]
These statements include \(\dim X_1=0\), for which full row rank is
the unrestricted rank-zero condition.
\end{srlemma}

\begin{proof}
Use the adapted bases
\((x,X_1,y)\) and \((x',W_1,y')\), and transport the split Fourier
matrix by the invertible Gram matrices. This does not change rank.
For a lower Fourier matrix \(C_1\), the Heisenberg invariances of
\(e_B\rho(a\,y\otimes x')\) force a possible upper Fourier matrix to
have block shape
\[
 C=\begin{pmatrix}
       0&0&\lambda\\
       0&C_1&b\\
       c&d&e
    \end{pmatrix}.                               \tag{T4}
 \label{eq:sr-theta-cone-cutoff-matrix}
\]
Indeed invariance under \(x\otimes W\) forces the first row, except
possibly its \(y'\) entry, to vanish; commuting the shift through this
invariance supplies that possible entry. Invariance under
\(X_1\otimes x'\) gives the zero first column in the middle block.
Intertwining the reduced Heisenberg action identifies the remaining
middle block as exactly \(C_1\). These are eigenvalue identities for
every local modulation, so the Fourier support is contained in (T4),
including for arbitrary linear combinations of test functions.

If \(C_1\) has a nonzero left kernel vector \(r\), its combination of
the middle rows has support only in the last column. If \(\lambda=0\)
the first row is already zero. If \(\lambda\ne0\), subtract the suitable
multiple of the first row from this combination. In either case
\(C\) has a nonzero left kernel vector. Thus a full upper row rank
forces full lower row rank. At \(a=0\) the commutator eigenvalue
\(\lambda\) is zero, so the first row itself is zero and the
evaluation vanishes. This proves (T2).

Finally \(e_A\) is, up to its fixed nonzero model scalar, the
composition of \(e_0\) with the evaluation invariant under
\(Ex\otimes W_1\). In the split Fourier model this invariance forces
the nonzero \(x\) row to vanish. A full row rank vector is therefore
killed, proving (T3). The argument only uses elementary row operations
and the actual modulation eigenvalues. In particular it does not
replace the split Fourier matrix by an unrelated coordinate matrix.
\end{proof}

\begin{srlemma}[The seminorm form of the convergence bounds]
\label{lem:sr-theta-integrated-seminorm}
Fix a finite dimensional space of finite test functions.
On the large side in its proved convergence range there are a
continuous infinite Schwartz seminorm \(p\) and an integrable
nonnegative function \(G\) on the first unitary quotient such that,
for \(h\) in any fixed compact subset of the second unitary group,
\[
 |\theta(f\otimes\phi;g,h)|\le p(f)G(g)
 \tag{T5}\label{eq:sr-theta-seminorm-domination}
\]
for each of a fixed finite basis of finite test functions.
The same assertion holds on the small side when that finite test
space lies in the actual full-rank cutoff.
The constants and seminorm may depend on this fixed test space,
the fixed model, and the compact set of \(h\).
They do not depend on \(f\).
\end{srlemma}

\begin{proof}
We give the upgrade from the coefficient estimates, since a mere
constant \(C_f\) for each individual \(f\) would not imply (T5).
In the large-side proof choose a single Schwartz order \(N\) greater
than the dimension and all exponents used in its lattice majorants.
Every occurrence of a test-function bound there is bounded by
\[
 P_N(f)=\max_{|\alpha|\le N}\sup_z
                 (1+\|z\|)^N|\partial^\alpha f(z)|.
\]
Finite linear coordinate changes, the finitely many finite
coefficient values, and the specified compact angular variables change
this to at most \(C P_M(f)\) for some fixed \(M,C\).
This follows directly by the chain rule; the matrices and all their
inverses range in a compact set. The root and lattice sums are then
exactly the previously proved integrable chamber majorant, multiplied
by \(C P_M(f)\).

For the small side, use the coefficient-square majorant in the
full-rank decay proof, keeping its test-function bound as \(P_N(f)\).
The compact-coset Fourier cancellation and the nonzero minor occur
only in the finite coefficients. They have a common level and common
finite coefficient bounds on a fixed finite basis. The remaining
Schwartz bounds, finite coordinate changes, and chamber summations
are linear inequalities in \(P_N(f)\). Thus their rapid chamber
majorant has coefficient \(C P_M(f)\), again with fixed \(M,C\).
Choose the decay order larger than the fixed Haar-growth exponent.
The resulting \(G\) is integrable; finite quotient volume handles the
remaining bounded chamber. This is precisely the same coefficient
argument as the preceding full-rank decay proof, with its constants
left in seminorm form.

For a compact family of \(h\), its finite action has finite orbit on
a fixed smooth finite vector. Its infinite action carries each
Schwartz seminorm into a uniformly bounded Schwartz seminorm.
In the present oscillator model this last assertion follows from
the explicit linear changes, Fourier transforms, and quadratic
phases on finitely many compact coordinate patches; their coefficients
are bounded. Taking the finite maximum proves the compact uniformity.

There is a small linear-algebra point when a family is known to lie
in the cutoff but its displayed finite tensor factors do not.
Let \(V\) be their finite dimensional span and let \(V^o\) be the
subspace satisfying the finite cutoff. Choose a projection
\(V\longrightarrow V^o\). At every infinite coordinate the finite
slice of such a family lies in \(V^o\), so applying this projection
does not change the family. Re-expand in a basis of \(V^o\).
Only finitely many constant linear combinations of infinite Schwartz
functions occur, and their seminorms are bounded by a fixed sum of
the original seminorms. This supplies exactly the cutoff basis needed
above.
\end{proof}

\begin{srlemma}[Joint integrability for unfolding a cone]
\label{lem:sr-theta-cone-joint-integrability}
Let \(F_l=e_A\rho(u_A(l))f\) on the large side, or
\(F_l=e_B\rho(u_B(l))f\) on the small side with \(f\in\mathcal S^o\).
Use the respective exponents
\[
 z_A=n-\frac d2,\qquad z_B=1+\frac d2-n,
\]
in the rank ranges in which they are at least \(2\), and in which
the large lower theta integral has the proved convergence range.
Then for every fixed lower second-group argument \(h_1\),
\[
 \int_{\mathbb A_E^\times}\int_{[U(X_1)]}
       |\theta(F_l;g_1,h_1)|\,dg_1\,|l|_E^{z_\epsilon}
                       \,d^\times l <\infty .
 \tag{T6}\label{eq:sr-theta-cone-joint-integrability}
\]
The assertion also holds after compact averaging in the upper first
unitary group. All exchanges of the cone integral with the lower
theta period are consequently valid.
\end{srlemma}

\begin{proof}
In the common partial Fourier coordinates an eliminated coordinate is
fixed to the nonzero corner parameter \(l_\infty\). Accordingly, for
every lower Schwartz seminorm \(q\) and every \(N\) there is an upper
Schwartz seminorm \(p\) such that the infinite factors satisfy
\[
 q(F_l)\le p(f)(1+\|l_\infty\|)^{-N}.
\]
This is the restriction estimate proved for the actual cone operator:
differentiate only in lower coordinates, bound their polynomial
weights by upper weights, and retain arbitrary extra weight in the
corner coordinate. At the finite places the argument in the finite
cone lemma leaves a finite dimensional space of lower test functions,
and forces the finite part of \(l\) into one fixed fractional ideal.
Near \(l=0\) these finite factors are locally constant.

On the \(B\) side (T2) puts every \(F_l\) in the actual lower cutoff.
Use the finite projection at the end of the preceding lemma to put
its entire finite test space in that cutoff. Now (T5), on either side,
bounds the inner integral by a fixed finite sum of the just displayed
Schwartz seminorms. The finite idele integral of the absolute
weight converges for \(z_\epsilon\ge2\): away from finitely many
places it is the Euler product
\(\prod_w(1-Nw^{-z_\epsilon})^{-1}\), and the finitely many remaining
tails are geometric series. At infinity write each complex coordinate
as \(r e^{i\theta}\). The weight and multiplicative Haar density give
\(r^{2z_\epsilon-1}\,dr\,d\theta\), integrable at zero since
\(z_\epsilon>0\). Choose \(N\) larger than the sum of all radial growth
exponents to obtain integrability at infinity.

The same estimates are uniform on the compact set being averaged.
This proves (T6), by an absolute majorant, and also proves convergence
in every Schwartz seminorm before applying any continuous lower theta
functional or lower theta period.
\end{proof}

\begin{srlemma}[The constant term of a single theta kernel]
\label{lem:sr-theta-kernel-constant-term}
Let \(N_W\) be the unipotent radical stabilizing \(Ex'\).
Give \([N_W]\) probability Haar measure. There are nonzero
\(c_0,c_1\), fixed by the model and evaluations before the finite
character packet, such that
\[
 \begin{split}
 \int_{[N_W]}\theta(f;g,uh)\,du
   ={}&c_0\Theta(e_0\omega(g,h)f)\\
     &+c_1\sum_{v\in X(E),\ v\ne0,\ \langle v,v\rangle=0}
       T_A(1,\omega(\gamma_vg,h)f).
 \end{split}
 \tag{T7}\label{eq:sr-theta-kernel-ct}
\]
Here \(\gamma_v\in U(X)(F)\) is chosen with \(\gamma_vv=x\).
If \(X\) is anisotropic, the second line is absent.
For each \(g,h\) the series is absolute.
\end{srlemma}

\begin{proof}
In the polarization \(W=Ex'\oplus W_1\oplus Ey'\), Poisson expansion
in the line coordinates writes the upper theta functional as the sum,
over \(v\in X(E)\), of the lower theta functional applied to
\(e_0\rho(v\otimes y')f\), with a fixed nonzero normalization.
The lattice Schwartz estimate in the first lemma proves absolute
convergence uniformly as the unipotent parameter ranges over its
compact adelic quotient. Thus it is legitimate to average this
expansion term by term.

Here is the orbit calculation behind that average. In the displayed
basis of \(W\), an element of \(N_W\) fixes \(x'\), sends
\[
 q\in W_1\ \longmapsto\ q-\langle q,z\rangle x',\qquad
 y'\ \longmapsto\ y'+z+
            \bigl(c-\tfrac12\langle z,z\rangle\bigr)x',
 \quad z\in W_1(\mathbb A_F),\ c\in\mathbb A_F .
\]
The skew Hermitian identities verify that these matrices preserve
the form. The center acts on the \(v\) term by the additive character
with coefficient \(\langle v,v\rangle\). Orthogonality on
\(F\backslash\mathbb A_F\) kills it unless this coefficient is zero.
At \(v=0\) the noncentral action also disappears, giving exactly the
first term in (T7). Its normalization is nonzero.

For nonzero isotropic \(v\), Witt decomposition over \(E/F\) supplies
\(\gamma_v\). A direct construction extends \(v\) to a hyperbolic
pair by choosing \(w\) with \(\langle v,w\rangle=1\) and subtracting
\(\langle w,w\rangle v/2\); Witt extension then identifies its
orthogonal complement with the fixed one. Rational invariance of
theta moves this term to the term with \(v=x\).
The remaining \(z\)-average is the average of the Heisenberg
translations in \(Ex\otimes W_1\). Fourier orthogonality in those
coordinates keeps exactly their zero modes. Its symplectic reduction
is \(X_1\otimes W_1\); in the common partial Fourier coordinates it is
the actual \(e_A\rho(x\otimes y')\) evaluation.
The fixed comparison of these coordinate evaluations gives \(c_1\ne0\).
This proves (T7). A different choice of \(\gamma_v\) changes the
expression by a rational stabilizer action, which has no effect on
the lower theta functional. No unitary splitting scalar is introduced
by either unipotent average: those matrices have determinant \(1\).
\end{proof}

\begin{srlemma}[Unfolding the two isotropic-vector sums]
\label{lem:sr-theta-isotropic-unfold}
In the convergence ranges of (T6), the respective isotropic-vector
sums have integrable absolute value on \([U(X)]\), and their integrals
are
\[
 \begin{split}
 \int_{[U(X)]}\sum_v T_A(1,\omega(\gamma_vg,h)f)\,dg
     &=k_A I_{X_1,W_1}(C_A\omega(1,h)f),\\
 \int_{[U(X)]}\sum_v T_B(1,\omega(\gamma_vg,h)f)\,dg
     &=k_B I_{X_1,W_1}(C_B\omega(1,h)f),
 \end{split}
 \tag{T8}\label{eq:sr-theta-isotropic-unfold}
\]
with fixed \(k_A,k_B\ne0\). The exponents in the cone operators are
exactly \(z_A,z_B\) in (T6). On the \(A\) line its character is
\(\chi_W^{-1}\), and on the \(B\) line its character is \(\chi_W\).
\end{srlemma}

\begin{proof}
The rational nonzero isotropic vectors form the homogeneous space
\(Q_x(F)\backslash U(X)(F)\), where
\(Q_x=N_X\rtimes U(X_1)\) is the stabilizer of the vector \(x\).
The integrand is invariant under \(Q_x(F)\) and under \(N_X(\mathbb A_F)\).
Unfold this rational sum, first with nonnegative absolute values.
Iwasawa decomposition writes the resulting quotient integral in
coordinates \(g=m_X(t,g_1)k\), with
\(g_1\in[U(X_1)]\), \(t\in\mathbb A_E^\times\), and \(k\in K_X\).

The Jacobian can be read off the same explicit unipotent matrices.
Conjugation by \(m_X(t,1)\) multiplies its \(n-2\) \(E\)-coordinates
by \(t\), and its central \(F\)-coordinate by \(t\bar t\).
Their additive module is \(|t|_E^{n-2}|t|_E=|t|_E^{n-1}\).
The quotient Haar density is consequently \(|t|_E^{1-n}\).
The volume of the rational unipotent quotient and the comparisons of
the chosen Haar measures give a positive geometric factor.

The actual mixed evaluation has first Levi multipliers
\[
 \begin{array}{c|c|c}
       &\hbox{multiplier of the evaluation}
                  &\hbox{corner after moving the Levi}\\ \hline
 A&\chi_W(t)|t|_E^{d/2-1}
                    &t^{-1}x\otimes y'\\
 B&\chi_W(t)|t|_E^{d/2}
                    &\bar t\,y\otimes x'.
 \end{array}                                      \tag{T9}
 \label{eq:sr-theta-evaluation-levi-table}
\]
To check the powers, the \(A\) eliminated space is
\(X\otimes x'+x\otimes W_1\): its determinant under the first Levi
is \(t^{d-2}\). The \(B\) space
\(x\otimes W+X_1\otimes x'\) has determinant \(t^d\).
The oscillator half-density gives half these powers.
The splitting multiplier is the prescribed \(\chi_W(t)\).
Conjugating the corner by the inverse first Levi gives the last
column. These computations use the same evaluation and same
oscillator action as the local constructions.

Multiplying by the Haar density, the \(A\) integrand has weight
\(\chi_W(t)|t|_E^{d/2-n}\); substitute \(l=t^{-1}\).
The \(B\) integrand has weight
\(\chi_W(t)|t|_E^{d/2-n+1}\), with \(l=t\).
This gives exactly the stated characters and exponents.
Joint integrability (T6), uniformly in \(k\), makes the unfolded
absolute integral finite. Tonelli therefore proves the asserted
summability almost everywhere and integrability of the rational
sums, and Fubini proves the complex equalities in (T8).
It also moves the cone operator through the lower theta period.
The geometric Haar and evaluation comparisons are nonzero and
independent of the finite character packet.
\end{proof}

\begin{srproposition}[The four actual constant-term producers]
\label{prop:sr-theta-four-constant-terms}
For a triple with upper ranks \(a,a'\), \(d=a+a'\), write
\[
 \chi_X=\beta,\quad \chi_W=\beta/\alpha
       \quad\hbox{on the \(V\) side},\qquad
 \chi_X=\alpha,\quad \chi_W=\alpha/\beta=\nu
       \quad\hbox{on the \(V'\) side}.
\]
Assume exactly the rank and signature ranges of the two adapted
codimension-two steps. Then the following hold with nonzero constants
fixed by the geometric models, evaluations, and Haar measures before
the finite packet.

On the large side the \(\beta|t|_E^{a/2}\) constant-term component is
the lower theta integral of \(e_0\omega(1,h)f\), times a fixed
nonzero constant. In the \(A\) step its
\(\alpha|t|_E^{a'/2}\) component is the lower large theta integral of
\[
       C_A\omega(1,h)f,\qquad z_A=(a-a')/2 .
 \tag{T10}\label{eq:sr-theta-large-cone-producer}
\]
On the small side with actual full-rank cutoff, in the \(B\) step
the \(\beta|t|_E^{a/2}\) component is the lower small theta integral of
\[
       C_B\omega(1,h)f',\qquad z_B=(a-a')/2+1 .
 \tag{T11}\label{eq:sr-theta-small-cone-producer}
\]
If \(V'\) is definite, the entire small constant term is its
\(\alpha|t|_E^{a'/2}\) component, equal to the lower theta integral of
\(e'_0\omega'(1,h)f'\), times a fixed nonzero constant.
These formulas hold on every specified lower Levi
\(m_W(1,h_1)\), with the lower theta integral evaluated at \(h_1\).
\end{srproposition}

\begin{proof}
First interchange the compact unipotent average with the first
theta period. On the large side (T5) gives an integrable majorant
uniform on this compact unipotent family. On the small side use the
same assertion on the actual cutoff, which is stable under the
second unitary action. Hence Fubini applies in both cases.
In (T7) the zero term has second Levi character and power
\(\chi_X(t)|t|_E^{n/2}\), by direct evaluation in the removed
\(W\) line. Changing the first quotient Haar measure to that of the
specified lower model contributes a fixed positive factor. This
already proves the large zero-term formula. If \(V'\) is definite
there are no nonzero rational isotropic vectors, so it proves the
entire small assertion as well. This includes rank zero.

On the large side (T8) applies to the nonzero orbit in (T7).
The rank inequality is precisely
\(2n=d+2z_A\) and
\(z_A\ge\min(r_0,s_0)+2\) for the lower signature \((r_0,s_0)\);
it puts the lower large integral in its proved convergence range.
The second Levi law of this cone is
\[
       (\chi_X/\chi_W)(t)|t|_E^{(d-n)/2}.
 \tag{T12}\label{eq:sr-theta-cone-second-law}
\]
It follows by the same mixed-coordinate computation as (T9), now
for the second Levi; scaling its line coordinate in the cone integral
gives the ratio of splitting characters. Since on this side
\(n=a\), the law is exactly
\(\alpha(t)|t|_E^{a'/2}\). This proves (T10).

For the small side it would be incorrect to unfold the \(A\) orbit
with absolute values in its divergent range. Instead apply (T1)
separately on every rational isotropic line in \(V'\).
Choose one representative \(v_L\) of each line \(L\), and let
\(\gamma_Lv_L=x\). The other nonzero vectors of that line are
\(\xi v_L\), \(\xi\in E^\times\). The first Levi formula, with
\(|\xi|_{\mathbb A_E}=1\) and trivial rational Hecke characters,
identifies their \(A\) terms with
\(T_A(\xi,\omega(\gamma_Lg,h)f')\).
For the \(B\) terms it gives the coefficient
\(\overline{\xi^{-1}}\); this is a permutation of \(E^\times\).
The zero \(B\) term vanishes by (T2).
Subtracting the zero terms from (T1) therefore gives, on each line,
\[
 \sum_{\xi\ne0}T_A(\xi,\omega(\gamma_Lg,h)f')
   =c_{AB}\sum_{\xi\ne0}
          T_B(\xi,\omega(\gamma_Lg,h)f')
                -T_A(0,\omega(\gamma_Lg,h)f').
 \tag{T13}\label{eq:sr-theta-small-paired-orbits}
\]
This identity is between absolutely convergent line series.

The \(B\) rational-vector sum is integrable by (T8), because its
lower functions retain the actual cutoff by (T2) and
\(z_B\ge2\). The full kernel constant term is integrable by the
compact-average majorant at the start of the proof. Thus the
difference of that kernel term and the \(B\) sum is integrable,
without asserting separate global absolute convergence of the
\(A\) orbit. By (T7) and (T13) this difference is
\[
 c_0\Theta(e_0\omega(g,h)f')
       -c_1\sum_L T_A(0,\omega(\gamma_Lg,h)f')
 \tag{T14}\label{eq:sr-theta-small-alpha-correction}
\]
wherever the \(B\) sum is summable. Both summands in (T14) have
second Levi law \(\chi_X(t)|t|_E^{n/2}\): the corner is now zero,
so evaluation has the zero-term law. The line series is summable
almost everywhere, by (T13) and the two rational-vector sums.
Consequently its integral, taken as the integrable combined
expression (T14), has that same law.

The remaining component is, by (T8), a fixed nonzero constant times
the lower small theta integral of \(C_B\omega(1,h)f'\).
Its law (T12) is now
\(\beta(t)|t|_E^{a/2}\), since \(n=a'\), whereas (T14) has law
\(\alpha(t)|t|_E^{a'/2}\). This proves (T11) and constructs the
required two-term splitting on the small side.

Finally these two Levi laws are distinct: \(a>a'\), and the
characters are unitary, so at an idele with \(|t|_E\ne1\) their
absolute values differ. A function having both laws is zero.
Subtracting any two two-term splittings shows their components
coincide. Hence the formulas just proved apply to the actual named
components of any specified splitting, not only to the splitting
used in this proof. Throughout, inserting a lower second Levi
\(m_W(1,h_1)\) intertwines the reduced oscillator action with the
specified lower model. The same Fubini arguments give the stated
formulas at every \(h_1\).
\end{proof}

\paragraph{Scope.}
\label{rem:sr-theta-ct-scope}
The small-side correction (T14) is essential. Omitting it, or
interchanging a divergent \(A\) orbit sum with an absolute integral,
would not prove the small cone formula. The construction above gives
the four constant-term identities and the exact cone exponents using
the actual full-rank cutoff and the actual evaluation maps. It still
does not infer descent of these maps through the upper matching
quotient merely from local scalarity. That descent and the global
orthogonality argument are the remaining steps in using these
identities for a rank recursion.

